% !PRINTED 286
\weberpage{286}

Nach dieser Bezeichnungsweise sind wir imstande, die Gruppe der Transformationsgleichung aus der der Teilungsgleichung sofort abzuleiten.

Wir setzen zunächst als Rationalitätsbereich den Inbegriff der rationalen Funktionen von $\varkappa^2$ mit \gesperrt{rationalen Zahlenkoeffizienten} fest.

Nach § 63, 3. besteht in diesem Rationalitätsbereich die Gruppe der Teilungsgleichung aus allen Substitutionen, durch die $\mu$, $\mu'$ in
\[
\begin{array}{c} \partial \mu - b \mu' \\ -c \mu + a \mu' \end{array}
\]
übergeführt werden, worin $a$, $b$, $c$, $\partial$ beliebige, nach dem Modul $p$ genommene ganze Zahlen sind, deren Determinante
\weq{9}{\Delta = a\partial - bc}
durch $p$ nicht teilbar ist, und die Anzahl aller dieser Substitutionen beträgt
\weq{10}{p(p-1)(p^2-1) \qquad [\S\ 63,\,(27)]}

Daraus ergibt sich aber nach der Bezeichnungsweise (4), (5) \gesperrt{die Gruppe der Transformationsgleichung als bestehend aus allen} durch das Symbol
\weq{11}{\left(z, \frac{\partial z - b}{-cz+a}\right) \ (\mathrm{mod}\ p)}
\gesperrt{ausgedrückten Vertauschungen.}

Wir bezeichnen eine Substitution $(z, z')$, wenn
\weq{12}{z \equiv \frac{c+\partial z'}{a+bz'} \ (\mathrm{mod}\ p)}
ist, ähnlich wie früher durch
\weq{13}{\begin{pmatrix} a, & b \\ c, & \partial \end{pmatrix},}
und erhalten für die Zusammensetzung zweier solcher Substitutionen, wenn
\[
z' \equiv \frac{c'+\partial' z''}{a'+b'z''} \ (\mathrm{mod}\ p)
\]
ist, die Regel:
\weq{14}{(z,z')(z'z'') \equiv \begin{pmatrix} a, & b \\ c, & \partial \end{pmatrix}\begin{pmatrix} a', & b' \\ c', & \partial' \end{pmatrix} \equiv \begin{pmatrix} aa'+bc', & ab'+b\partial' \\ ca'+\partial c', & cb'+\partial\partial' \end{pmatrix} \ (\mathrm{mod}\ p),}
%CONT>>
